\documentclass{article}
\usepackage{amsmath}
\usepackage{amssymb}
\usepackage{amsthm}
\usepackage{physics}


\newtheorem*{conjecture}{Conjecture}
\newtheorem*{theorem}{Theorem}
\newtheorem*{lemma}{Lemma}


\makeatletter
\newcommand*{\rom}[1]{\expandafter\@slowromancap\romannumeral #1@}
\makeatother


\begin{document}

\section*{A useful fact for absolute value}

\begin{lemma}
	For any \(x \in \mathbb{R}\), \(|x| \geq x\)
\end{lemma}
\begin{proof}
	Suppose \(x \geq 0\). Then from the definition of the absolute value: \(|x| = x \geq x\). Suppose now that \(x < 0\). We see:
	\begin{equation*}
		|x| = -x > 0 > x
	\end{equation*}
	So \(-x > x\) which implies \(-x \geq x\).
\end{proof}

\end{document}